% Shared content for the Bitcoin Austria style demo. Included verbatim by both
% demo.tex (16:9 — the standard) and demo-43.tex (4:3), so the two aspect-ratio
% builds differ ONLY in slide geometry. That is exactly what we want to compare:
% how the theme lays out the very same slides in each format.
%
% The style package is one level up; it self-locates its fonts/ and assets/.
\usepackage{../bitcoin-austria}

\title{Bitcoin Austria}
\subtitle{Wenn dein Geld wirklich nur dir gehört.}
\author{Bitcoin Austria}
\date{\today}

\begin{document}

\begin{frame}[plain,noframenumbering]
  \titlepage
\end{frame}

\section{Wer wir sind}

% Corner image: a cropped, captioned meetup photo in the top-right corner.
% Placed below the headline band so it never collides with the logo/title.
% The image floats ON TOP of the slide and does not reflow text, so we keep the
% body in a narrower left column that stays clear of (and reads beside) the photo.
\cornerimage[width=0.20\paperwidth, viewport=150 90 760 600, clip,
             border, caption=Wiener Meetup]{meetup}
\begin{frame}{Über uns}
  % Left-aligned narrow column so the bullets read beside the corner photo,
  % never under it (the corner image is an overlay and does not reflow text).
  \noindent\begin{minipage}[t]{0.68\textwidth}
    \begin{itemize}
      \item Die \textbf{älteste Bitcoin-Organisation der Welt} — gegründet
        2011 aus den ersten Wiener Meetups
      \item Über \textbf{200 Meetups und Events} mit tausenden Teilnehmenden
      \item Gemeinnützig und unabhängig — finanziert ausschließlich durch
        \textbf{Bitcoin-Spenden}
      \item Eine zivilgesellschaftliche Stimme für unabhängige
        Bitcoin-Expertise in Österreich
    \end{itemize}
  \end{minipage}

  \vfill
  \begin{block}{Unser Auftrag}
    Bitcoin Austria vermittelt unabhängiges Wissen, damit du informierte
    Entscheidungen über dein Geld triffst – gemeinnützig, seit 2011.
  \end{block}
\end{frame}

\fillerslide[Drei Prinzipien, die nicht verhandelbar sind]{Unsere Werte}

\section{Unsere Werte}

\begin{frame}{Drei Prinzipien}
  \begin{itemize}
    \item \textbf{Bitcoin Only} — ein dezentrales Netzwerk, kein spekulatives
      Krypto-Asset
    \item \textbf{Privatsphäre} — nicht verhandelbar; ein klares Nein zu
      programmierbaren CBDCs
    \item \textbf{Wissen statt Marketing} — Erklärungen statt Anlageberatung
  \end{itemize}

  \vfill
  \begin{alertblock}{Wofür wir stehen}
    Wirtschaftliche Selbstbestimmung — im Geist der Österreichischen
    Schule.\footnote{Wirtschaftswissenschaftliche Denkrichtung um Menger, Mises
      und Hayek — betont freie Märkte und individuelle Selbstbestimmung.}
  \end{alertblock}
\end{frame}

\comparisonslide{Bitcoin ist nicht „Krypto“}%
  {Bitcoin}{%
    \begin{itemize}
      \item Dezentrales Netzwerk, kein Emittent
      \item Feste, vorhersehbare Geldmenge
      \item 100\,\% Open Source
      \item Läuft seit 2011 ununterbrochen
    \end{itemize}}%
  {„Krypto“}{%
    \begin{itemize}
      \item Marketing und Spekulation\footnote{Mit „Krypto“ sind hier die
        tausenden spekulativen Token rund um Bitcoin gemeint — nicht Bitcoin
        selbst.}
      \item Oft zentral kontrolliert
      \item Beliebig vermehrbar
    \end{itemize}}%
  {Warum wir unterscheiden}{%
    Wir erklären Bitcoin — und grenzen es klar von spekulativen
    Krypto-Projekten ab.}

% Corner image on a theme-generated layout: bottom-right, full (uncropped) photo.
\cornerimage[corner=bottom-right, width=0.16\paperwidth, border]{meetup}
\begin{frame}{Bitcoin vs. Fiat-Geldsystem}
  \centering
  \comparisontable{Bitcoin}{Fiat-Geldsystem}{%
    Verwahrung    & du selbst         & Banken \\
    Zugang        & erlaubnislos      & Konto nötig \\
    Verfügbarkeit & rund um die Uhr   & Bankzeiten \\
    Zensur        & keine             & möglich \\
    Geldmenge     & fix (21 Mio.)     & beliebig vermehrbar \\
  }
\end{frame}

% Simplest call: bare width, default top-right corner.
\cornerimage[0.17\paperwidth]{meetup}
\begin{frame}{Mitmachen}
  \begin{itemize}
    \item Bildungsinhalte zu Technik, Ökonomie und Gesellschaft
    \item Meetups, öffentliche Events und Pressearbeit
    \item Gespräche mit politischen Entscheidungsträgerinnen und -trägern
    \item Förderung von Open Source und einer offenen Community
  \end{itemize}

  \vfill
  \begin{center}
    „Damit mehr Menschen Bitcoin verstehen — nicht, damit jemand daran verdient.“

    \medskip
    \large \textbf{bitcoin-austria.at}
  \end{center}
\end{frame}

\end{document}
